\documentclass[11pt]{article}

\usepackage[utf8]{inputenc}
\usepackage[T1]{fontenc}
\usepackage[margin=1in]{geometry}
\usepackage{amsmath,amssymb}
\usepackage{booktabs}
\usepackage{tabularx}
\usepackage{xcolor}
\usepackage{setspace}
\usepackage[colorlinks=true,linkcolor=blue,urlcolor=blue]{hyperref}

\setstretch{1.12}

\definecolor{revcol}{RGB}{150,20,20}
\definecolor{defcol}{RGB}{20,70,150}
\definecolor{sevhigh}{RGB}{150,20,20}
\definecolor{sevmed}{RGB}{190,120,0}
\definecolor{sevlow}{RGB}{20,120,40}
\definecolor{concede}{RGB}{120,60,140}

% concern block commands ------------------------------------------------------
\newcommand{\concern}[2]{%
  \par\medskip\noindent{\large\textbf{#1}}\hfill{\small #2}\par
  \vspace{2pt}\hrule\vspace{4pt}}
\newcommand{\Rev}[1]{\noindent\textcolor{revcol}{\textbf{$\blacktriangleright$ Reviewer.}}\ \textit{#1}\par\smallskip}
\newcommand{\Def}[1]{\noindent\textcolor{defcol}{\textbf{$\blacktriangleleft$ Author response.}}\ #1\par}
\newcommand{\Concede}[1]{\smallskip\noindent\textcolor{concede}{\textbf{Concession / change made.}}\ #1\par}
\newcommand{\sevH}{\textcolor{sevhigh}{\textbf{[Severity: High]}}}
\newcommand{\sevM}{\textcolor{sevmed}{\textbf{[Severity: Medium]}}}
\newcommand{\sevL}{\textcolor{sevlow}{\textbf{[Severity: Low]}}}

\title{\textbf{Adversarial Review and Author Defense}\\[4pt]
\large of \emph{``Memorization or Recognition? A Gold-Standard Mahabharata NER\\
Corpus and a Feature-Engineering Framework that Out-Generalizes a\\
Pretrained Transformer''}}
\author{Prepared as a self-critique: a skeptical reviewer and a defending author,\\
arguing over the same evidence.}
\date{\today}

\begin{document}
\maketitle

\begin{abstract}
\noindent This companion document stress-tests the paper. A deliberately hostile
reviewer raises the objections a strong program committee would raise --- data
leakage, an unfair transformer baseline, the relevance of classical models,
annotation quality, sample size, and statistical power --- and the authors
respond with the evidence already in the manuscript and its supporting run
artifacts. Where an objection is partly right, the authors concede it explicitly
and point to the change made; a defense that never concedes is not a defense but
a denial. The central thesis survives scrutiny: the headline claim is deliberately
the \emph{hardest} number in the paper (span-$F_1$ on \emph{unseen} entity
surfaces), not the easy one, and the evaluation is constructed precisely to keep
memorization from masquerading as recognition.
\end{abstract}

\section{What the paper claims (so we agree on the target)}
The paper makes three claims. (C-I) It releases a human gold-annotated
\emph{Mahabharata} NER corpus: $67{,}482$ sentences, $119{,}927$ confirmed entity
mentions, nine domain-specific types, across all nineteen \emph{parvas}, built by
human correction of a transformer's pre-labels ($13.2\%$ relabeled). (C-II) Over a
shared, explicit handcrafted feature framework $\phi(w_i,S)$ --- with a circular
name-gazetteer deliberately excluded --- classical models and a CRF are trained
and compared, on the gold labels, against \texttt{en\_core\_web\_trf} used
zero-shot. (C-III) The honest, leakage-controlled result is that on \emph{unseen}
entity surfaces the CRF ($0.787$) and Logistic Regression ($0.792$) beat the
transformer ($0.466$) by $\sim\!0.32$ span-$F_1$, while a causal ablation
attributes that generalization to morphology and context rather than token
identity. The overall $\sim\!0.98$ span-$F_1$ is reported as a
\emph{memorization ceiling}, not the result.

Everything below attacks and then defends these three claims.

\section{Reviewer's overall assessment (the hostile framing)}
\textcolor{revcol}{\itshape
``The paper reports classical models scoring $0.98$ span-$F_1$ on NER --- a figure
that in this subfield is a red flag for leakage. It then claims to `beat
transformers,' but the transformer is never trained on the data. The gold set is
built by correcting the very model it is later compared against, which invites
circularity. Results are on a $35$k-sentence sample, not the full corpus, and the
key `unseen' claim rests on only $587$ spans with no confidence intervals. I am not
yet convinced the headline is real rather than an artifact of evaluation design.''}

\medskip
\noindent The authors take each of these seriously below. The short version: the
$0.98$ is explicitly \emph{not} the claim; the transformer comparison is honestly
scoped as ``off-the-shelf'' and the fair arena (unseen) is the one where the
transformer is \emph{not} disadvantaged by the split; the annotation protocol's
anchoring risk is real and is bounded; and the sample-size and power concerns are
valid limitations that do not touch the internal validity of the comparison.

% =====================================================================
\section{Part A --- Evaluation validity and leakage}
% =====================================================================

\concern{A1. ``0.98 span-$F_1$ is impossibly high --- this is leakage.''}{\sevH}
\Rev{Classical token classifiers reaching $0.98$ span-$F_1$ on NER is not
credible without leakage. Modern neural systems on CoNLL-2003 sit around
$0.90$--$0.94$. A Linear SVM at $0.982$ should not exist. Something has leaked
from test into training.}
\Def{The number is real and its mechanism is fully explained, and it is
\emph{not} test-into-train leakage. Two facts. First, we verified the pipeline:
duplicate sentences are removed \emph{before} the $70/15/15$ split (so no sentence
appears in two folds), and every label-dependent statistic --- frequency bins and
the entity prior $P(y\mid w)$ --- is fitted on the \emph{training fold only}.
There is no path by which a test label reaches a training feature. Second, the
reason the number is high is stated in the paper: in a natural sentence-level
split of a \emph{narrative} corpus, entity surfaces recur heavily. Concretely,
$94.0\%$ of test-fold entity mentions ($9{,}139$ of $9{,}726$) have a surface
string that already occurs as a gold entity in training. The $\phi^{prior}$ and
$\phi^{core}$ (token-identity) features then let \emph{any} learner recover the
training-majority label for those surfaces almost perfectly. That is
\emph{memorization within a legitimate split}, not leakage across the split ---
and it is precisely why we refuse to report $0.98$ as the result.}
\Def{The CoNLL comparison is apt but cuts the other way: CoNLL-2003 test entities
overlap far less with training, and its metric is not dominated by a recurring
closed cast of characters. The correct comparison to CoNLL is our \emph{unseen}
number ($\sim\!0.79$), not our overall number.}

\concern{A2. ``Then the $0.98$ is still deceptive --- you are burying it.''}{\sevM}
\Rev{Even if it is not leakage, foregrounding a memorization-inflated $0.98$
anywhere is misleading readers who skim tables.}
\Def{We agree that a bare $0.98$ would mislead, which is why the manuscript takes
active countermeasures rather than passive ones. The leaderboard's caption
literally reads ``\emph{the overall span-$F_1$ column is a memorization-inflated
ceiling},'' the body paragraph immediately following it sets the number in bold as
\textbf{not} the paper's claim, and the very next subsection replaces it with the
seen/unseen decomposition. The abstract leads with the unseen gap, not the
overall score. A reader would have to ignore the caption, the paragraph, and the
abstract to be misled.}

\concern{A3. ``The seen/unseen split is the only thing keeping this honest --- is it correct?''}{\sevH}
\Rev{Your entire honesty argument rests on the definition of `seen.' If that is
computed loosely --- say, by token rather than by entity surface, or including
validation --- the claim collapses.}
\Def{The definition is the strict one. A test entity span is ``seen'' iff its
exact (lower-cased) surface string appears as a \emph{gold entity} anywhere in the
\emph{training} fold; otherwise it is ``unseen.'' Validation is not counted as
training for this purpose. Scoring is done separately on the two partitions with
identical span-matching. The resulting supports ($n_{\text{seen}}=9{,}139$,
$n_{\text{unseen}}=587$) sum exactly to the test entity total ($9{,}726$), which
is a consistency check that the partition is exhaustive and non-overlapping. On
the unseen partition, $\phi^{prior}$ and $\phi^{core}$ can contribute nothing by
construction, so any score there is attributable to the generalizing families.}

\concern{A4. ``The `prior' feature is a memorized lexicon in disguise.''}{\sevM}
\Rev{$P(y\mid w)$ is just a soft gazetteer. Including it is the same sin you
accuse others of --- your model is a lookup table with extra steps.}
\Def{This is exactly right for \emph{seen} words, and we neither hide it nor
depend on it. The paper isolates the prior in two ways. (i) The seen/unseen split:
on unseen surfaces the prior returns the uniform distribution and cannot help, yet
the model still reaches $0.79$. (ii) The causal ablation: removing the prior from
the full model costs only $0.980\!\rightarrow\!0.971$, and the prior \emph{alone}
reaches $0.553$ --- respectable, but far below the full model, and clearly a
memorization channel we label as such. The genuinely irreplaceable family is not
the prior but \emph{context}: removing it collapses the model to $0.838$. Our
contribution is not that the prior is good; it is that we can \emph{separate} the
prior's memorization from the features' recognition and report each honestly.}

\concern{A5. ``The unseen number is only $587$ spans --- statistically meaningless.''}{\sevM}
\Rev{$587$ spans, no confidence intervals, no significance test. A $+0.32$ gap
could be noise, and one lucky fold could produce it.}
\Def{We accept that $587$ is modest and that CIs would strengthen the claim.
Three points bound the concern. First, the effect is not marginal: $0.79$ vs
$0.47$ is a $32$-point gap, and even a conservative binomial interval on a few
hundred trials does not span that distance. Second, the pattern is coherent across
independent evidence --- the per-type table (transformer at $0.00$ on
\textsc{Concept}/\textsc{Weapon}/\textsc{Event}, $0.10$ on \textsc{Dynasty}) and
the correction analysis both predict exactly this failure, so the unseen gap is
not an isolated data point. Third, a validation-scale replication ($2{,}000$
sentences) produced the same ordering. }
\Concede{We add bootstrap confidence intervals over test sentences for the
seen/unseen scores, and a larger-corpus rerun to grow the unseen support, to
future work, and we now state the $587$ figure explicitly in the table caption
rather than only in the text.}

\concern{A6. ``Unseen-$F_1$ is inflated by the easy \textsc{Person} majority.''}{\sevM}
\Rev{Most unseen spans are \textsc{Person}, and any capitalized proper noun
defaults to \textsc{Person} correctly. Your $0.79$ measures the easy case.}
\Def{Correct, and stated in the paper. Unseen-$F_1$ is a \emph{micro} average in
which \textsc{Person} dominates, so $0.79$ should be read as ``recognizes an unseen
name as an entity of a plausible type,'' not as perfect fine-grained typing; the
harder unseen non-\textsc{Person} types are lower. Crucially, this caveat
\emph{weakens our absolute number but not our comparative claim}: the transformer
faces the same easy \textsc{Person} majority and still scores only $0.47$, because
it over-tags \textsc{Person} and cannot place the rest. The gap is what matters,
and the majority helps both sides.}
\Concede{We add an unseen \emph{macro}-$F_1$ (averaged over attested types on the
unseen slice) as the fairer summary in the camera-ready.}

% =====================================================================
\section{Part B --- Baselines and fairness}
% =====================================================================

\concern{B1. ``Comparing a trained model to a zero-shot one is not a fair fight.''}{\sevH}
\Rev{Your $\phi$ models are trained on the gold training fold; the transformer is
never fine-tuned. Of course the trained model wins. The comparison is rigged and
the title's `out-generalizes a pretrained transformer' is an overclaim.}
\Def{This is the most important objection and we have scoped the claim to survive
it. The comparison is deliberately an \emph{off-the-shelf} setting: it answers the
practitioner's real question --- ``can I just run the best pretrained NER on this
epic, or must I engineer features?'' --- and the answer is that off-the-shelf NER
fails ($0.68$ overall, $0.00$ on several native types). We never claim classical
models beat transformers \emph{in general}; the manuscript says so in three places
(baseline subsection, seen/unseen qualifications, discussion). On the
\emph{unseen} slice the framing is actually \emph{more} favorable to the
transformer than it looks: the transformer's seen and unseen scores are nearly
equal ($0.698$ vs $0.466$), proving it gets \emph{no} benefit from the training
vocabulary, so the split does not advantage it or disadvantage it --- the unseen
column is the one arena where ``trained vs zero-shot'' is least of an issue,
because neither side is memorizing.}
\Concede{A transformer \emph{fine-tuned} on the gold training fold is the correct
stronger baseline and is named as the single most important item of future work.
We soften no result, but we explicitly frame the current transformer as a lower
bound on what transformers can do here, not an upper bound.}

\concern{B2. ``No neural sequence baseline (BiLSTM-CRF, fine-tuned BERT).''}{\sevM}
\Rev{A serious NER paper in 2026 needs at least one \emph{trained} neural baseline
to contextualize the classical numbers.}
\Def{We agree this is the main empirical gap and do not pretend otherwise; the
earlier draft's claim of BiLSTM-CRF/mBERT baselines was removed precisely because
we did not run them, rather than left in as vaporware. What the paper \emph{does}
establish is self-contained and still novel: the gold corpus, the leakage-controlled
protocol, and the demonstration that cheap, interpretable features generalize to
unseen names far better than the strongest \emph{off-the-shelf} transformer. Adding
a fine-tuned neural baseline changes the competitive landscape but not any of the
three contributions.}

\concern{B3. ``The OntoNotes$\rightarrow$schema mapping cripples the transformer unfairly.''}{\sevM}
\Rev{You mapped the transformer's OntoNotes labels into your nine types. Types
like \textsc{Dynasty}, \textsc{Concept}, \textsc{Weapon} have no OntoNotes
equivalent, so the transformer scores $0.00$ by construction. That is your mapping's
fault, not the model's.}
\Def{The $0.00$ scores are a \emph{finding about domain transfer}, not a mapping
artifact, and we report them as such. The mapping is generous, not adversarial:
we route \textsc{Org}$\rightarrow$\textsc{Dynasty} and
\textsc{Work\_of\_Art}$\rightarrow$\textsc{Sacred\_Text} to give the transformer
every reasonable chance. The point stands that a general-domain label space simply
does not contain the epic's ontology --- which is exactly why a gold corpus with a
domain schema was needed in the first place. We also do not rest the headline on
these zeros: even on \textsc{Person}, the one type squarely in its wheelhouse, the
transformer ($0.84$) trails the CRF ($0.98$).}

% =====================================================================
\section{Part C --- Data and annotation}
% =====================================================================

\concern{C1. ``Correcting the model's own labels invites anchoring bias.''}{\sevH}
\Rev{Annotators started from \texttt{en\_core\_web\_lg}'s output. Humans anchor to
a suggested label and under-correct. Your gold is therefore biased \emph{toward}
the model you later evaluate --- circular and self-serving.}
\Def{Anchoring is a genuine risk of correct-a-baseline annotation and we do not
dismiss it. Three mitigations bound it. (i) The correction rate is high and
structured, not cosmetic: annotators changed $13.2\%$ of mentions
($16{,}310/123{,}994$), deleted $4{,}067$ spurious mentions outright, and
\emph{reduced} \textsc{Person} from the model's $96{,}175$ to $83{,}129$ ---
i.e.\ they corrected against the model's dominant bias, which is the opposite of
anchoring to it. (ii) If anchoring were severe, it would inflate the transformer's
apparent score by making the gold agree with it; instead the transformer scores
only $0.68$, so any residual anchoring biases \emph{against} our own headline, not
for it. (iii) The pre-tagging model (\texttt{lg}) is \emph{different} from the
evaluated model (\texttt{trf}), so there is no identity between annotator prior and
test system.}
\Concede{We add an explicit annotation-bias discussion and recommend a
double-annotated, adjudicated subset with inter-annotator agreement (e.g.\ Cohen's
$\kappa$) in the corpus release, which the current version lacks.}

\concern{C2. ``No inter-annotator agreement is reported.''}{\sevM}
\Rev{Without IAA the `gold' label is one person's opinion. Ambiguous cases
(\textsc{Dynasty} vs \textsc{Person} for `Kaurava') have no measured reliability.}
\Def{Fair. The corpus's reliability on precisely these hard boundaries
(\textsc{Dynasty}$\leftrightarrow$\textsc{Person}, \textsc{Concept}$\leftrightarrow$\textsc{Person})
is unquantified in this version, and those boundaries are exactly where the CRF's
residual errors concentrate --- so the model's ``errors'' there may partly reflect
label uncertainty rather than model failure.}
\Concede{We flag IAA on a re-annotated sample as required for the corpus paper and
temper any per-type claim on the lowest-agreement types accordingly.}

\concern{C3. ``Footnote-digit noise (`Panchala35') means the text is dirty.''}{\sevL}
\Rev{Surfaces like \texttt{Kurus2} in your error analysis show OCR/footnote
contamination survived preprocessing and is corrupting both training and scoring.}
\Def{Correct in the small, and quantified: such artifacts are a named,
minority error mode in the CRF analysis (a handful of the $145$ false negatives),
not a pervasive corruption. Preprocessing strips the common cases; the residue is
a real limitation of working from a scanned classical text, and we report it in
the error analysis rather than sweeping it up. It bounds the achievable ceiling
but does not change the seen/unseen comparison, which is symmetric across models.}

% =====================================================================
\section{Part D --- Methods, framing, and scope}
% =====================================================================

\concern{D1. ``Why classical ML at all in the transformer era?''}{\sevM}
\Rev{This reads as nostalgia. Feature engineering lost. Why should a 2026 reader
care about SVMs and CRFs?}
\Def{Because the paper's question is not ``what wins a leaderboard'' but ``what is
the cheapest, most \emph{interpretable} thing that actually recognizes unseen
entities in a low-resource, specialized domain, and \emph{which} linguistic
signals carry that ability.'' A CRF over $\phi$ trains in minutes on a CPU, needs
no GPU, and --- through the causal ablation --- tells a Sanskrit scholar that
\emph{morphology} and \emph{context}, not memorized names, drive recognition. That
is a scientific result about the language, not merely a system. Transformers can
be added on top; they do not answer this question, and off-the-shelf they do not
even clear $0.68$.}

\concern{D2. ``The SVM result specifically looks anomalous.''}{\sevL}
\Rev{A Linear SVM at $0.982$, second only to the CRF, is suspicious for NLP.}
\Def{It is not an anomaly and it is not the top result. The ranking is
CRF $0.983 >$ LinearSVM $0.982 \approx$ LogReg $0.981$ --- a tight cluster of
\emph{linear} models, exactly as decades of statistical NLP predict for sparse,
high-dimensional, largely one-hot features; linear SVMs and CRFs were the standard
strong classical sequence labelers long before neural nets. The SVM is a member of
the leading cluster, not an outlier, and it rides the same seen-surface
memorization as every other model in the overall column (see A1). Its \emph{unseen}
behavior is not even reported as headline because we retain span-aligned
predictions only for the CRF and LR; the SVM is not part of the generalization
claim.}

\concern{D3. ``SHAP/LIME are advertised but the SHAP figure is missing.''}{\sevL}
\Rev{The interpretability section promises SHAP and LIME, but the primary evidence
is an ablation, and the SHAP artifact appears incomplete.}
\Def{This is deliberate, and defensible on methodological grounds independent of
any missing artifact. Raw $\sum|\text{SHAP}|$ is biased by family cardinality ---
the context family alone contributes tens of thousands of sparse columns and can
dominate the sum by count rather than importance. We therefore designate the
\emph{causal} family ablation (retrain with each family added/removed) as the
primary evidence and use SHAP only as qualitative corroboration. The scientific
claim about which families matter rests on the ablation, which is complete and
reported in full.}

\concern{D4. ``$35$k sample, not the full $67$k corpus --- results may not hold.''}{\sevM}
\Rev{You ran on a $35{,}000$-sentence sample for memory reasons. The headline may
be a sampling artifact; report the full corpus.}
\Def{The sample is a uniform draw over the deduplicated corpus and is split
$70/15/15$ identically, so it is representative of the whole in composition (the
$\sim\!91\%$ \texttt{O}-rate and per-type proportions match the full-corpus
figures). We also built and validated a streaming, checkpointed pipeline that fits
the full sparse matrix without holding it in memory, so scaling is an engineering
run, not a redesign.}
\Concede{We agree the full-corpus run should appear in the final version, and we
present the $35$k results as representative rather than final, stating the sample
size prominently in the experimental setup.}

\concern{D5. ``Does any of this generalize beyond one epic?''}{\sevL}
\Rev{This is a single-corpus study. The `finding' may be idiosyncratic to the
\emph{Mahabharata}.}
\Def{The corpus is specific; the \emph{method} is not. The transferable recipe ---
release human gold labels, exclude circular lookup features from headline models,
and always report seen and unseen separately --- applies to any specialized or
historical domain where entity surfaces recur and general-domain taggers
under-perform. We position the epic as a demanding case study of a general
evaluation pathology, and name extension to further Sanskritic and Indo-European
texts as future work.}

% =====================================================================
\section{Scorecard}
% =====================================================================
Table~\ref{tab:scorecard} summarizes every concern, its severity, whether the
defense holds, and what (if anything) changed in response. The pattern is the
honest one: the internal-validity attacks (leakage, seen/unseen definition,
circular prior) are \emph{refuted} with evidence; the external-validity and
completeness attacks (fine-tuned baseline, IAA, full-corpus run, CIs) are
\emph{conceded} as limitations and future work. No concession touches the three
core contributions.

\begin{table}[!htbp]
\centering
\small
\renewcommand{\arraystretch}{1.25}
\caption{Reviewer concerns, severity, and disposition.}
\label{tab:scorecard}
\begin{tabularx}{\textwidth}{@{}l l l X@{}}
\toprule
\textbf{ID} & \textbf{Concern} & \textbf{Severity} & \textbf{Disposition} \\
\midrule
A1 & $0.98$ is leakage                 & High   & \textbf{Refuted} --- dedup-before-split, train-only stats; $0.98$ is a memorization ceiling, not leakage \\
A2 & $0.98$ is deceptive               & Medium & \textbf{Refuted} --- flagged in caption, body, and abstract as \emph{not} the claim \\
A3 & seen/unseen split wrong           & High   & \textbf{Refuted} --- strict surface-in-train definition; supports sum to test total \\
A4 & prior is a hidden lexicon         & Medium & \textbf{Refuted} --- isolated by split + ablation; context, not prior, is irreplaceable \\
A5 & $n{=}587$ underpowered            & Medium & \textbf{Partly conceded} --- add bootstrap CIs + larger run; effect size is large \\
A6 & unseen inflated by \textsc{Person}& Medium & \textbf{Partly conceded} --- add unseen macro-$F_1$; comparative gap unaffected \\
B1 & trained vs zero-shot unfair       & High   & \textbf{Refuted + scoped} --- honest off-the-shelf framing; unseen arena neutral to the split \\
B2 & no neural baseline                & Medium & \textbf{Conceded} --- fine-tuned transformer is the key future baseline \\
B3 & OntoNotes mapping cripples trf    & Medium & \textbf{Refuted} --- mapping is generous; zeros are a domain-transfer finding \\
C1 & anchoring bias in annotation      & High   & \textbf{Refuted + concession} --- corrections oppose the model's bias; add IAA subset \\
C2 & no inter-annotator agreement      & Medium & \textbf{Conceded} --- add $\kappa$ on re-annotated sample \\
C3 & footnote/OCR noise                & Low    & \textbf{Refuted} --- minority, reported error mode; symmetric across models \\
D1 & why classical ML                  & Medium & \textbf{Refuted} --- interpretability + the scientific ``which features'' question \\
D2 & SVM anomalous                     & Low    & \textbf{Refuted} --- linear cluster behind CRF; historical norm \\
D3 & SHAP incomplete                   & Low    & \textbf{Refuted} --- ablation is the principled primary evidence \\
D4 & $35$k sample                      & Medium & \textbf{Partly conceded} --- representative sample; full run to camera-ready \\
D5 & single-corpus generality          & Low    & \textbf{Refuted} --- method/recipe generalizes; corpus is a case study \\
\bottomrule
\end{tabularx}
\end{table}

% =====================================================================
\section{Closing statement for the defense}
% =====================================================================
The reviewer's sharpest instinct --- that a $0.98$ NER score signals leakage --- is
the right instinct applied to the wrong number. The paper's design anticipates
exactly this reflex and defuses it: the $0.98$ is dissected into a $94\%$ seen
partition (memorization, disclosed) and a $6\%$ unseen partition (recognition,
headline), the circular gazetteer is removed from every headline model, and all
label statistics are fitted on training data behind a deduplicated split. What
remains after every legitimate deduction is a modest but real claim: on names it
has never seen, an interpretable feature-based model recognizes epic entities
substantially better than the best off-the-shelf transformer, and the ability
comes from morphology and context rather than memorized surface forms.

The concessions --- a fine-tuned neural baseline, inter-annotator agreement,
confidence intervals, and the full-corpus run --- are real and are the right
next steps. None of them can convert a memorization ceiling into the headline, and
none of them weakens the two artifacts the community can use today: a
human gold-annotated \emph{Mahabharata} NER corpus, and an evaluation protocol
that refuses to let a recurring cast of characters be mistaken for recognition.
We rest the defense there.

\end{document}
